% !Mode:: "TeX:UTF-8"
\chapter{绪论}
\label{cha:chap01}

\section{研究背景与意义}

细胞是生命活动的基本结构与功能单元。细胞群体在基因表达、代谢状态、形态结构、迁移能力和力学表型等方面具有显著差异，群体平均结果难以完整反映单个细胞的状态分布。单细胞测序、质谱流式细胞术、基于液滴条形码的单细胞转录组测序和高通量成像等技术推动了细胞分析由群体统计向单细胞分辨发展\cite{altschuler2010heterogeneity,trapnell2015singlecell,shapiro2013singlecell,bendall2011singlecellmass,macosko2015dropseq,klein2015dropletbarcoding,heath2016singlecelltools,regev2017humancellatlas}。单细胞力学表型与分子信息和形态信息相互补充，可用于描述细胞在外部力学环境中的状态和响应。

细胞力学行为由细胞膜、膜皮质层、细胞骨架、胞质和细胞核等结构共同决定。细胞骨架承担力的传递与重构，细胞膜和膜皮质层维持边界与整体形态，胞质和细胞核参与时间相关的变形与恢复\cite{kasza2007cellmaterial,fletcher2010cellmechanics,wang2009mechanotransduction,moeendarbary2013poroelastic}。基质刚度能够改变细胞黏附、迁移和分化行为\cite{discher2005tissue,engler2006matrix}。肿瘤细胞、血细胞和免疫细胞的弹性、黏弹性及变形能力也会随细胞类型和状态发生变化，因而细胞力学表型可作为细胞状态表征中的一类物理指标\cite{suresh2007cancer,dicarlo2012mechanical,toepfner2018morphorheological}。

细胞弹性模量、黏性系数、黏弹性特征时间、膜张力和整体变形能力通常不能由静态图像直接获得。这些参数需要在外部载荷和边界条件已知的情况下，根据力与位移之间的关系、形态变化、时间响应或恢复过程进行估计。原子力显微镜、微吸管吸吮、光学拉伸和磁镊等方法能够获得局部或整体细胞力学响应，但通常依赖逐细胞定位、局部加载或较复杂的实验操作。高通量细胞力学测量则要求在连续处理大量细胞的同时保持加载条件、观测量和参数解释之间的一致性\cite{darling2015high,dicarlo2012mechanical,urbanska2020comparison}。

细胞形态与细胞力学参数之间存在联系，但二者并不等同。静态形态受到细胞类型、细胞周期、渗透状态和成像姿态等因素影响，外形相近的细胞可能具有不同的材料性质。外部载荷作用后的动态形态能够记录变形建立、峰值变形和形态恢复过程，其中变形幅值与细胞刚度相关，响应滞后和恢复速度包含黏弹性信息，质心运动还受到流体载荷和几何阻力影响\cite{mietke2015rtdc,mokbel2017rtdc,dudani2016viscoelastic,gong2024prediction}。因此，动态形态为细胞力学参数测量提供了比单一静态轮廓更丰富的信息。

微流控通道可利用压力驱动流动、通道收缩、剪切流和伸展流对细胞施加可控载荷，并在细胞运动过程中连续记录轮廓和位置。细胞轮廓、质心位置、通过时间和运动速度共同描述细胞受力和变形的时空过程。实时变形细胞术和定量变形细胞术等方法提高了细胞形态和力学表型的测量通量\cite{otto2015rtdc,nyberg2017qdc,rosendahl2018rtdc}。微流控条件下的细胞形态同时受到流体载荷、通道几何、细胞尺寸和材料参数影响，因而需要结合力学模型解释形态指标与力学参数之间的关系\cite{mietke2015rtdc,mokbel2017rtdc}。

实际微流控成像受到帧率、视场范围、曝光条件和图像分割误差影响。完整形态序列可能被压缩为少量观测时刻，局部轮廓也可能受到边界噪声干扰。单帧形态描述特定位置处的变形状态，完整序列还包含变形建立、峰值变形和恢复过程。多时刻轮廓、细胞位置和时间信息已被用于弹性模量预测\cite{gong2024prediction}。在稀疏观测条件下，还需要研究缺失形态的重构以及重构结果中的力学信息保持问题。

载细胞液滴（cell-laden droplet）下降接触固壁构成另一类瞬态流体加载过程。液滴接触固壁后经历铺展、最大铺展、回缩、反弹或沉积，液滴外界面和接触线运动由惯性、黏性、表面张力和壁面润湿性共同控制\cite{rein1993impact,yarin2006impact,josserand2016drop,snoeijer2013moving}。当细胞位于液滴内部时，液滴接触引起的压力和速度场变化通过液体传递至细胞表面，细胞的变形和运动又会影响局部流动。载细胞液滴可同时提供液滴铺展、动态接触角、细胞轮廓和质心运动等动态信息。Tasoglu等采用复合液滴模型研究载细胞液滴沉积过程及内部细胞变形\cite{tasoglu2010compound}。复合液滴撞击研究也表明，内部组分会改变外部液滴的铺展、回缩和界面运动\cite{blanken2021compoundperspective}。

微通道细胞变形和载细胞液滴下降接触均利用外部流体载荷驱动细胞变形，并以动态形态作为力学参数反演的观测信息。微通道的几何边界和压力条件相对稳定，细胞响应主要由通道约束、流动条件和细胞材料参数共同决定。载细胞液滴下降接触还包含气液自由界面、表面张力和壁面润湿等作用，外部液滴形态与内部细胞形态共同反映瞬态载荷和细胞力学响应\cite{tasoglu2010compound,blanken2021compoundperspective,freund2014bloodcells}。

由形态响应估计力学参数可采用模型标定、迭代优化和数据驱动预测等方法。模型标定使用解析关系或预先计算的响应曲线，参数含义较为明确。迭代优化通过反复求解正问题使模拟形态接近观测形态，计算成本随参数维数和正问题复杂度增加。数据驱动方法在预先建立的数据库上学习形态响应与力学参数之间的非线性关系，可提高参数预测效率，其可靠性取决于数据库覆盖范围、数据划分方式和观测信息\cite{raue2009identifiability,roberts2017crossvalidation,vabalas2019limited}。

\section{细胞力学特性的测量与形态响应反演}

细胞力学测量通常包括外部加载、响应观测和参数计算三个环节。传统方法通过探针、吸管、光场或磁场施加局部或整体载荷，微流控方法通过连续流动和微尺度通道诱导细胞变形。不同方法在加载范围、测量通量、观测时间尺度和参数计算方法方面存在差异。形态响应反演则以细胞轮廓、位移和时间变化为输入，在给定力学模型或数据驱动模型下估计细胞力学参数。

\subsection{传统单细胞力学测量方法}

原子力显微镜（atomic force microscopy，AFM）通过悬臂梁探针与细胞表面接触，记录加载力与压入深度之间的关系。Binnig等提出AFM，为微纳尺度表面形貌和相互作用力测量建立了仪器基础\cite{binnig1986afm}。细胞力学实验通常将压入曲线与Hertz模型或其修正形式结合，计算细胞表观弹性模量。探针形状、压入深度、加载速率、细胞黏附状态和基底效应均会影响模量估计结果\cite{gavara2016afm}。AFM适合测量细胞局部刚度及其空间差异，也可用于比较药物处理或病理状态变化前后的细胞力学性质\cite{rotsch2000cytoskeletal,cross2007nanomechanical}。

微吸管吸吮通过负压将细胞局部吸入微吸管，并记录吸入长度随压力和时间的变化。薄壳、液滴和黏弹性模型可用于计算膜张力、皮质张力、表观黏度和整体变形参数。Hochmuth系统总结了活细胞微吸管吸吮中的力学模型和实验处理方法\cite{hochmuth2000micropipette}。Evans和Yeung根据血液粒细胞的吸入过程计算表观黏度和皮质张力，体现了该方法对整体黏弹性响应的表征能力\cite{evans1989granulocytes}。

光镊和光学拉伸器利用光场对细胞或微粒施加载荷。单束梯度力光阱为微尺度非接触操控提供了基础\cite{ashkin1986singlebeam}。光学拉伸器采用两束相向激光对悬浮细胞施加整体拉伸，通过长轴和短轴变化及恢复过程描述细胞变形能力\cite{guck2001optical}。光学变形性已被用于区分不同细胞状态\cite{guck2005opticaldeformability}。磁镊和粒子示踪微流变方法则通过磁性微粒或示踪粒子测量局部黏弹性\cite{bausch1999local,waigh2005microrheology}。

传统单细胞力学测量具有较明确的加载边界和较强的参数解释能力。AFM侧重局部压入，微吸管吸吮侧重整体吸入和膜皮质层响应，光学方法侧重非接触整体拉伸，磁镊和微流变侧重局部黏弹性。不同方法测得的弹性模量和黏性参数对应不同加载尺度、时间尺度和模型假设，比较结果时需要考虑具体测量条件\cite{dicarlo2012mechanical,urbanska2020comparison}。

\subsection{基于微流控的高通量细胞力学测量方法}

变形细胞术（deformability cytometry，DC）利用微流控通道中的压力、流动和几何结构使细胞发生变形，并通过高速成像或其他检测手段获得细胞力学表型。按照细胞的主要加载方式，常用微流控变形细胞术可分为收缩通道变形细胞术、剪切流变形细胞术和伸展流变形细胞术\cite{urbanska2020comparison}。

收缩通道变形细胞术（constriction-based deformability cytometry，cDC）利用尺寸小于或接近细胞直径的狭窄通道，使细胞在压力驱动下发生受限变形。细胞通过时间、进入时间、堵塞概率和变形轮廓可用于表征细胞变形能力。Shelby等建立感染疟原虫红细胞通过毛细通道的微流控模型，展示了细胞变形能力与微循环阻塞之间的关系\cite{shelby2003capillary}。Byun等利用变形和表面摩擦区分不同癌细胞的力学表型\cite{byun2013deformability}。微流控棘轮结构也可依据细胞变形能力进行分选\cite{park2016ratchets}。由于细胞需要与通道壁接触，壁面摩擦、黏附和堵塞可能影响测量结果。

伸展流变形细胞术（extensional flow deformability cytometry，xDC）利用交叉流或流动聚焦结构形成伸展流场，使悬浮细胞在短时间内发生拉伸变形。Gossett等利用惯性微流控中的水动力拉伸实现了大样本单细胞力学表型测量\cite{gossett2012hydrodynamic}。Dudani等在微流控伸展流中测量细胞随时间变化的变形，并利用加载和恢复过程表征细胞黏弹性\cite{dudani2016viscoelastic}。xDC减少了细胞与固壁直接接触，但变形结果仍受到流场分布、细胞尺寸、初始位置和成像时刻影响。

剪切流变形细胞术（shear flow deformability cytometry，sDC）利用微通道内的水动力应力使悬浮细胞发生变形。实时变形细胞术（real-time deformability cytometry，RT-DC）是代表性的sDC方法。Otto等提出RT-DC，实现了连续流动条件下的高通量单细胞力学表型测量\cite{otto2015rtdc}。Mietke等结合低雷诺数流动和线弹性理论，建立了由细胞尺寸和变形量估计细胞刚度的方法\cite{mietke2015rtdc}。实时荧光与变形细胞术（real-time fluorescence and deformability cytometry，RT-FDC）进一步将荧光信号与细胞形态和力学表型同步测量\cite{rosendahl2018rtdc}。

图~\ref{fig:rtdc_mietke2015}给出了RT-DC中的微通道结构、观测位置和随细胞运动的坐标描述。

\begin{figure}[htbp]
  \centering
  \includegraphics[width=0.82\textwidth]{figures/chap01/fig_rtdc_mietke2015.png}
  \caption[实时变形细胞术中的通道结构与受限流动示意图]{实时变形细胞术中的通道结构与受限流动示意图。图（a）给出了微通道几何、细胞变形区域和图像采集位置。图（b）给出了随细胞运动的坐标系及细胞周围的流动结构。引自文献\cite{mietke2015rtdc}。}
  \label{fig:rtdc_mietke2015}
\end{figure}

定量变形细胞术（quantitative deformability cytometry，q-DC）通过标定流体应力和时间相关应变获得细胞力学参数\cite{nyberg2017qdc}。不同DC方法在加载方式、通量、壁面接触、可测参数和模型依赖性方面具有不同适用范围\cite{urbanska2020comparison}。面积、圆度偏差、长宽比和通过时间等量首先属于细胞力学表型。将这些观测量转换为弹性模量或黏弹性参数，还需要给出流体载荷、通道几何、细胞模型和参数标定方法。

\subsection{微通道细胞变形与力学参数反演}

细胞在变截面微通道中运动时，依次经历进入收缩段、持续变形、通过狭窄区域和离开约束后的恢复过程。入口附近的轮廓主要反映初始形态和流动扰动，收缩段内的轮廓记录压力、黏性剪切和几何限制共同作用下的变形，出口后的轮廓和速度反映细胞恢复与整体运动。完整过程中的轮廓、位置和时间信息构成由形态响应估计材料参数的数据基础\cite{mietke2015rtdc,mokbel2017rtdc,gong2024prediction}。

微通道形态的力学解释依赖流固耦合模型。细胞受到流体压力、黏性剪切、壁面约束和材料恢复力共同作用，其变形不仅由弹性模量决定，还与细胞尺寸、泊松比、黏性参数和通道几何有关。Mokbel等采用数值方法研究RT-DC中细胞力学性质的提取，比较不同材料模型和皮质层参数对细胞形态的影响\cite{mokbel2017rtdc}。图~\ref{fig:rtdc_mokbel2017}显示，不同本构模型和参数设置会产生不同的细胞轮廓，三维形态中的局部凹陷也可能难以由单一侧视图像完整观察。

\begin{figure}[htbp]
  \centering
  \includegraphics[width=0.90\textwidth]{figures/chap01/fig_rtdc_mokbel2017.png}
  \caption[不同力学模型和参数条件下的RT-DC细胞形态预测结果]{不同力学模型和参数条件下的RT-DC细胞形态预测结果。上排比较不同细胞力学模型和参数条件下得到的轮廓。下排给出三维细胞后端形态及其侧向观测差异。引自文献\cite{mokbel2017rtdc}。}
  \label{fig:rtdc_mokbel2017}
\end{figure}

变截面微通道使细胞在有限距离内经历连续变化的流体载荷和几何约束，不同材料参数对应的轮廓差异可在多个位置被记录。Xie和Hu分析了胶囊通过局部收缩区域时的边界变形和通过行为，说明膜力学参数与通道形状共同决定变形过程\cite{xie2021capsule}。此类数值模拟建立材料参数与形态响应之间的正向关系，参数反演则根据观测轮廓估计相应的材料参数。

材料大变形、黏弹性、壁面接触和多阶段运动会增加解析反演的难度。数据驱动模型可由数值样本学习轮廓、位置、时间与力学参数之间的非线性关系。Gong等将多个时刻的细胞轮廓、时间和空间位置信息输入卷积神经网络，用于预测弹性模量和材料本构类别\cite{gong2024prediction}。该研究说明，多时刻形态包含单帧几何指标未覆盖的过程信息。

完整轮廓序列的获取受到成像帧率、曝光时间、视场范围和分割误差限制。观测帧数减少时，峰值变形或恢复阶段可能缺失。边界噪声还会改变局部曲率、周长和轮廓高频特征。点坐标表示能够直接保存轮廓局部位置，傅里叶轮廓表示则可将闭合边界转换为有限阶系数，以较低维度描述整体形态和局部变化\cite{kuhl1982elliptic}。形态、质心位置和速度的联合分析有助于同时描述细胞变形和整体运动\cite{caruana1997multitask}。

\section{载细胞液滴下降接触过程及其数值模拟研究}

\subsection{液滴撞击与壁面润湿}

液滴撞击固壁是典型的自由界面流动问题。液滴接触前的运动主要由初始速度和重力决定，接触后惯性促使液体沿壁面铺展，黏性耗散和表面张力抑制铺展并推动回缩。根据流动条件和壁面润湿性，液滴可能发生沉积、回缩、反弹或飞溅\cite{rein1993impact,yarin2006impact,josserand2016drop}。雷诺数、韦伯数、奥内佐格数和接触角常用于表征惯性、黏性、表面张力和壁面润湿作用。

最大铺展因子是液滴撞击研究中的常用指标。Pasandideh-Fard等结合实验和能量分析研究液滴最大铺展，说明初始动能、表面能和黏性耗散共同决定最大铺展直径\cite{pasandidehfard1996capillary}。Mao等建立液滴铺展和回弹模型，分析接触角和能量耗散对撞击结果的影响\cite{mao1997spread}。Richard等研究超疏水表面上的液滴接触时间，表明惯性和毛细作用决定回弹时间尺度\cite{richard2002contact}。高速撞击下还可能出现液片抬升和飞溅\cite{riboux2014experiments}。

壁面润湿性通过平衡接触角和动态接触角影响接触线运动。动态接触角随接触线速度、液体黏度、表面张力和壁面状态变化，铺展和回缩阶段通常具有不同的角度变化规律\cite{deGennes1985wetting,bonn2009wetting,snoeijer2013moving}。因此，铺展因子、液滴高度、接触线位置和动态接触角可共同描述液滴接触固壁后的外界面演化。

\begin{figure}[htbp]
  \centering
  \includegraphics[width=0.90\textwidth]{figures/chap01/fig_drop_impact_pasandideh1996.png}
  \caption[水滴撞击固体表面的实验图像与数值结果]{水滴撞击固体表面的实验图像与数值结果。图中展示了液滴接触、铺展和回缩过程中的典型界面形态。引自文献\cite{pasandidehfard1996capillary}。}
  \label{fig:drop_impact_pasandideh1996}
\end{figure}

\subsection{载细胞液滴中的细胞力学响应}

载细胞液滴下降接触固壁时，液滴外界面演化、内部流动和细胞变形同时发生。液滴接触壁面后，轴向运动受到抑制，液体向径向流动，内部压力和速度场随铺展过程快速变化。细胞受到液体压力和黏性作用后发生变形和运动，其弹性和黏弹性又会影响形态恢复过程。Tasoglu等建立载细胞液滴沉积模型，分析液滴沉积条件和细胞变形之间的关系\cite{tasoglu2010compound}。相关研究表明，液滴尺寸、撞击速度、流体性质和基底条件均会影响载细胞液滴沉积过程中的细胞响应。

与均匀液滴相比，载细胞液滴内部存在可变形细胞，外部液滴形态和内部细胞形态具有不同的力学含义。液滴铺展主要反映惯性、黏性、表面张力和壁面润湿作用，细胞变形和恢复则与局部流体载荷、细胞尺寸和材料参数相关。内部细胞还会改变局部流道和液体运动，使液滴外界面与细胞形态之间产生相互影响\cite{blanken2021compoundperspective,freund2014bloodcells}。

载细胞液滴可提供液滴外轮廓和细胞轮廓两类动态形态。外部响应包括铺展因子、液滴高度、接触线位置和动态接触角，内部响应包括细胞长宽比、圆度偏差、质心位移、速度和恢复过程。材料参数首先影响细胞自身的变形与恢复，下降速度、表面张力和壁面接触角则直接改变液滴流动和外界面演化。不同参数可能对同一响应量产生相近影响，因此需要结合多个时刻和多类形态量分析参数敏感性和可辨识性\cite{raue2009identifiability}。

\subsection{载细胞液滴的数值模拟方法}

载细胞液滴下降接触的数值模拟需要同时处理气液自由界面、移动接触线和细胞变形。体积分数法通过网格单元中的体积分数表示不同流体占比，具有较好的质量守恒性质\cite{hirt1981vof}。水平集方法采用连续函数隐式表示界面，便于计算界面法向和曲率\cite{osher1988fronts,sussman1994levelset}。前沿追踪方法显式维护界面网格，可保持较清晰的界面位置\cite{unverdi1992fronttracking}。

相场法通过有限厚度的弥散界面描述气液界面，并将自由能、表面张力和界面演化统一写入控制方程\cite{anderson1998diffuseinterface,jacqmin1999phasefield,badalassi2003phasefield,yue2004diffuseinterface}。壁面自由能或相应边界条件可用于给定平衡接触角。细胞边界可采用膜、壳或连续体模型描述，具体选择决定材料参数的物理含义和可输出的力学量\cite{freund2014bloodcells,barthesbiesel2016motion}。

流体与细胞之间的相互作用可采用浸没边界法或任意拉格朗日欧拉方法处理。浸没边界法在固定流体网格上通过离散核函数传递结构力\cite{peskin2002immersed}。任意拉格朗日欧拉方法使流体网格随流固界面运动，可在界面处直接满足速度连续和牵引平衡\cite{donea2004ale}。载细胞液滴模型还需同时处理相场界面和流固界面，两类界面分别描述气液边界和细胞边界。

数值模型的可靠性需要通过基准算例、网格无关性和时间步长无关性进行检验。单液滴铺展、最大直径和接触角可用于验证自由界面和润湿模型，可变形结构在流体载荷下的响应可用于验证材料模型和流固耦合方法。网格和时间步分析应覆盖液滴铺展、细胞变形、质心运动以及应力等主要结果\cite{roache1998verification}。

\section{人工智能在形态重构与力学参数反演中的应用}

形态响应数据可表示为图像、轮廓点、几何指标和时间序列。形态重构以稀疏或含噪观测为输入，输出完整轮廓或连续形态过程。参数反演以一个或多个时刻的形态响应为输入，输出弹性模量、黏弹性参数或其他待识别参数。重构任务关注轮廓误差和几何量误差，反演任务关注参数预测误差及不同参数之间的可辨识性。

\subsection{数据驱动的形态表征与参数预测}

早期数据驱动材料建模主要采用前馈神经网络拟合应力与应变之间的关系或材料状态变量之间的映射。Ghaboussi等将神经网络用于材料行为表示，展示了由数据学习非线性材料响应的基本形式\cite{ghaboussi1991materialnn}。数据驱动计算力学进一步将材料数据与平衡条件、相容条件和边界条件联系起来\cite{kirchdoerfer2016datadriven}。细胞形态反演在此基础上还需要处理高维边界和时间序列。

多层感知机（multilayer perceptron，MLP）适合处理固定长度向量，可将多个观测时刻的轮廓坐标或几何指标拼接后进行回归。卷积神经网络（convolutional neural network，CNN）利用局部卷积核提取相邻轮廓点或图像区域的空间特征，适用于规则采样轮廓和二维图像\cite{lecun2015deeplearning,krizhevsky2012imagenet,he2016resnet}。细胞轮廓属于闭合曲线，使用CNN时需要保持轮廓点的排列方向、起始点和首尾邻接关系。

变分自编码器（variational autoencoder，VAE）将高维输入编码为潜变量分布，并由潜变量重构原始数据\cite{kingma2014vae}。该方法可用于学习轮廓的低维表示和重构缺失形态。Transformer采用注意力机制描述序列中不同位置之间的关系\cite{vaswani2017attention}。在稀疏形态重构中，不同观测时刻可作为序列输入，时间或空间位置可通过位置编码输入模型。

轮廓输入可分为点坐标、参数曲线和几何指标。点坐标保留局部边界细节，但维度较高并依赖轮廓点对应关系。面积、圆度偏差和长宽比等几何指标维度较低、物理含义明确，但不能完整描述局部边界。傅里叶轮廓表示可在保留主要形态特征的同时降低轮廓维度\cite{kuhl1982elliptic}。不同输入表示应根据重构精度、参数预测精度和计算复杂度进行比较。

\subsection{模型评价与可靠性分析}

形态重构需要同时评价轮廓局部误差和整体几何误差。点坐标误差反映对应轮廓点之间的距离，面积、周长和圆度偏差反映整体形态是否保持。重构序列还可输入参数预测模型，检验重构误差是否影响弹性模量等力学参数的估计。

参数反演需要根据预测对象选择绝对误差、相对误差和决定系数等评价指标。多参数任务还应分别报告各参数误差，避免不同量纲和数值范围掩盖单个参数的预测表现。参数敏感性并不等同于参数可辨识性，不同参数产生相近形态时，需要结合多时刻信息、内外形态和误差分布分析参数混淆\cite{raue2009identifiability}。

数据划分应以完整数值过程为基本单位。同一过程中的不同时间帧、稀疏组合和噪声版本不能同时进入训练集和测试集，否则会造成数据泄漏。随机划分主要评价给定参数范围内的插值能力，未见参数组合和参数范围边缘应单独测试\cite{roberts2017crossvalidation,vabalas2019limited}。重复训练和不同随机种子可用于评价结果对模型初始化和样本划分的敏感性\cite{ovadia2019uncertainty}。

% ------------------------------------------------------------------
% 以下两节按当前工作安排冻结。待第三章至第五章完成后统一修订。
% ------------------------------------------------------------------
\section{现有研究不足与本文研究思路}

现有细胞力学测量、微流控变形分析、复合液滴动力学和数据驱动参数预测已形成较完整的方法体系。传统测量提供明确载荷和参数解释，微流控方法提高形态数据获取通量，数值模拟给出材料参数与动态响应之间的正向关系，机器学习用于表示复杂形态与参数之间的非线性映射。本文研究涉及的主要不足集中在观测完整性、形态表示、参数可辨识性和模型可靠性四个方面。

第一，稀疏形态观测中的力学信息保持缺少统一评价。现有微流控参数估计通常使用单帧几何指标、指定时刻图像或较完整的形态序列。实际成像可能只覆盖少量位置和时刻，轮廓分割还会引入局部边界误差。稀疏到完整的轮廓恢复首先是几何重构问题，同时也是力学信息恢复问题。平均点坐标误差较低的轮廓可能改变面积、周长、曲率或关键变形阶段，从而影响弹性模量预测。现有评价较少同时覆盖局部轮廓、整体几何、质心运动和下游参数识别。

本文将微通道细胞变形过程划分为多个空间和时间阶段，以少量观测轮廓恢复完整过程。点坐标表示用于评价局部边界误差，面积、圆度和长宽比用于评价整体形态，固定参数预测模型用于评价重构结果中的弹性模量信息。数据按照完整数值过程划分，避免同一轨迹的不同时间帧跨越训练集和测试集。含噪输入与无噪输入分别用于检验几何稳定性和力学信息保持。

第二，高维轮廓数据的低维表示和动态信息利用仍不充分。离散点坐标能够完整记录边界，但输入维度高，对轮廓起点、点序和局部噪声较敏感。面积、圆度和变形指数等低维指标具有明确含义，却无法描述局部边界和形态方向。完整时间序列还包含细胞位置和速度信息，单一轮廓表示难以同时描述形态变化和整体运动。

本文采用傅里叶轮廓表示压缩闭合边界，并由傅里叶系数计算形态派生量。质心位置和速度作为独立动态变量进入联合预测任务。不同傅里叶阶数用于比较低维表示能力和局部细节保持，模型结构、残差分支和损失项通过消融实验评价。该部分与点坐标重构形成递进关系：前者检验局部边界从稀疏观测到完整过程的恢复，后者检验低维形态、运动和弹性模量的联合预测。

第三，含细胞复合液滴中的多参数响应存在混淆。外部液滴铺展同时受下降速度、表面张力、黏度和接触角影响，内部细胞变形同时受流体载荷、材料参数、尺寸和空间位置影响。某一参数变化能够改变响应，不代表该响应可唯一确定该参数。弹性模量变化和下降速度变化可能产生相似的最大变形，表面张力变化和接触角变化可能产生相似的外部铺展。不同参数还可能在早期和晚期表现出不同敏感性。

本文从标准工况出发，分别比较材料、流动、界面和几何参数变化下的内部形态、质心运动、应力、应变能以及外部铺展和动态润湿响应。分析重点由单条曲线的变化方向扩展到响应幅度、峰值时刻、恢复过程和内外响应之间的关系。参数--响应敏感性用于筛选具有区分能力的观测量，响应相关性和相近形态用于识别参数混淆。第四章给出正问题和可辨识性基础，第五章利用相同参数体系检验由动态形态反演参数的稳定性。

第四，小样本条件下的多参数反演、泛化和不确定性证据不足。多参数数据库的规模受数值计算成本限制，高维图像或多帧轮廓模型容易在有限独立过程中产生过拟合。随机划分主要评价已覆盖参数域内的插值，不能替代未见参数组合和范围边缘测试。模型平均误差也无法反映参数混淆区和高误差样本。

本文以完整计算过程为独立样本，明确训练集、验证集和测试集的参数分布。单参数和多参数任务分别评价，图像、轮廓和几何指标作为不同输入形式进行比较。帧数组合、内部与外部形态、模型结构和损失项通过消融实验分析。噪声输入、参数组合留出和范围边缘测试用于评价稳健性，重复划分或多随机种子结果用于估计训练不确定性。预测结论限定于本文数值模型、参数范围和观测方式。

第五，数值形态与实际观测之间存在模型和观测域差异。微通道与复合液滴数值模型采用确定的材料关系、边界条件和几何设置，真实细胞具有结构异质性和个体差异。数值轮廓可直接由计算边界提取，实验轮廓还受到成像分辨率、投影方向、光学模糊和分割算法影响。数值数据库适合检验方法、响应敏感性和参数可辨识性，其结果不能直接替代实验标定。

本文在数值域内统一形态提取、坐标对齐、采样和归一化规则，并通过人工扰动检验观测误差影响。第三章重点处理稀疏帧和轮廓噪声，第四章明确各响应量的后处理定义，第五章比较不同观测输入对参数预测的影响。实验迁移、真实细胞结构和跨设备泛化作为当前工作的适用边界，不在数值结果中作超出证据的定量推断。

基于上述问题，本文按照“形态恢复--响应分析--参数反演”的顺序组织研究。微通道研究以受控几何和压力加载为条件，检验少量轮廓能否恢复完整形态，并由重构结果估计弹性模量。傅里叶表示将局部点坐标问题扩展为低维形态、质心运动和弹性模量的联合预测。含细胞液滴研究以多相界面和瞬态流固耦合为条件，分析多类参数对内外动态形态的影响，并检验不同参数的可辨识性。动态形态参数预测在此基础上比较单参数、多参数、不同输入和不同测试分布下的误差。

\section{本文主要研究内容与章节安排}

本文研究外部流体载荷作用下细胞动态形态与力学参数之间的关系。微通道场景采用压力驱动和几何受限加载，含细胞液滴场景采用下降接触产生的瞬态惯性、自由界面和润湿加载。两类场景均由数值模拟生成形态数据，研究内容分别面向稀疏形态恢复、低维动态表示、多参数响应分析和力学参数预测。

本文的第一项研究内容为微通道稀疏轮廓重构与弹性模量信息保持。变截面微通道数值模型生成细胞通过收缩区域的动态轮廓。完整过程按照空间位置和时间顺序进行对齐，少量观测时刻构成稀疏输入。MLP、VAE 和 Transformer 等模型用于恢复未观测时刻的轮廓。结果评价包括边界点误差、面积、周长、圆度、长宽比和噪声稳定性。重构轮廓随后输入固定弹性模量预测模型，用于判断几何接近是否同时对应材料参数信息保持。

第二项研究内容为基于傅里叶轮廓表示的形态、质心运动与弹性模量联合预测。细胞闭合边界转换为有限阶傅里叶系数，质心位置和速度作为独立运动变量。联合模型由稀疏观测预测完整形态过程、质心速度和弹性模量。傅里叶阶数分析用于确定低维表示与局部细节之间的平衡，模型比较和消融实验用于评价网络结构、残差分支和损失项的作用。不同通道位置的误差和几何派生量用于分析模型在进入、压缩和恢复阶段的表现。

第三项研究内容为含细胞复合液滴下降接触过程的动力学及参数敏感性。数值模型同时描述外部液滴界面、壁面润湿、内部流体和细胞变形。标准工况用于划分接触、铺展、最大变形和回缩恢复等阶段。弹性模量、黏弹性特征时间、表面张力、下降速度、内部尺寸和壁面接触角等参数分别改变，内部细胞形态、质心运动、应力、应变能以及外部液滴铺展和动态接触角同步比较。综合分析给出不同响应对参数的敏感阶段和参数混淆关系。

第四项研究内容为基于复合液滴动态形态的力学参数预测。第四章生成的内外形态序列构成参数数据库，细胞轮廓、液滴轮廓、几何指标和多帧图像作为候选输入。单参数预测用于评价受控条件下的映射精度，多参数联合预测用于检验参数同时变化时的可辨识性。不同输入形式、帧数、模型和损失设置进行对比，噪声、参数组合留出和范围边缘测试用于评价稳健性与泛化边界。预测结果按照物理单位报告，并结合误差分布和高误差样本分析参数混淆。

全文共分为六章，各章内容如下。

第一章为绪论。本章介绍细胞力学参数与形态响应的研究背景，梳理传统单细胞力学测量、微流控高通量形态测量、微通道参数反演、含细胞复合液滴下降接触、数值模拟以及人工智能参数预测的研究进展，归纳稀疏观测、形态表示、参数可辨识性、小样本和数值--实验域差异等问题，并给出本文的研究内容。

第二章为数学模型与研究方法。本章给出两类研究场景采用的流体控制方程、自由界面模型、细胞材料模型、流固耦合条件和 ALE 有限元方法，统一定义形态与运动指标、数据预处理方式、机器学习任务和评价指标。具体几何、参数范围、网格、时间步和网络设置分别在对应研究章节中给出。

第三章为基于稀疏形态响应的微通道细胞力学参数反演方法。本章建立变截面微通道细胞变形数据库，研究点坐标条件下的稀疏轮廓重构和噪声稳定性，并通过弹性模量预测评价重构轮廓的力学信息。傅里叶轮廓表示用于形态、质心运动和弹性模量的联合预测，模型结果通过阶数分析、基线比较、结构消融和区域误差进行评价。

第四章为含细胞复合液滴下降接触过程的数值模拟与力学响应分析。本章建立多相流动--可变形结构耦合模型，完成自由界面、流固耦合、网格和时间步相关验证。标准工况给出内外形态的阶段性演化，不同材料、流动、界面和几何参数的影响通过内部变形、质心运动、应力、应变能、铺展因子和动态接触角进行比较。综合敏感性和可辨识性分析用于确定第五章的输入响应和目标参数。

第五章为基于含细胞复合液滴动态形态的力学参数智能预测。本章定义单参数和多参数反演任务，建立多帧图像、轮廓和几何特征输入，给出模型结构、训练设置和数据划分。预测结果包括逐参数误差、输入和帧数消融、噪声稳定性、未见参数组合、参数混淆和不确定性分析。

第六章为全文总结与展望。本章按照稀疏形态恢复、动态响应分析和力学参数反演三个问题总结主要结果，归纳本文工作的创新内容，说明数值模型、样本规模、实验迁移和多参数可辨识性的限制，并讨论实验验证、三维模型、真实细胞结构、域适应和不确定性量化等后续研究方向。
